
\documentclass{article}
\usepackage{colortbl}
\usepackage{makecell}
\usepackage{multirow}
\usepackage{supertabular}

\begin{document}

\newcounter{utterance}

\centering \large Interaction Transcript for game `clean\_up', experiment `2\_hard\_5obj\_it', episode 1 with nemotron{-}3{-}ultra.
\vspace{24pt}

{ \footnotesize  \setcounter{utterance}{1}
\setlength{\tabcolsep}{0pt}
\begin{supertabular}{c@{$\;$}|p{.15\linewidth}@{}p{.15\linewidth}p{.15\linewidth}p{.15\linewidth}p{.15\linewidth}p{.15\linewidth}}
    \# & $\;$A & \multicolumn{4}{c}{Game Master} & $\;\:$B\\
    \hline

    \theutterance \stepcounter{utterance}  
    & & \multicolumn{4}{p{0.6\linewidth}}{
        \cellcolor[rgb]{0.9,0.9,0.9}{
            \makecell[{{p{\linewidth}}}]{
                \texttt{\tiny{[P1$\langle$GM]}}
                \texttt{Sono il tuo maestro di gioco e stai giocando un gioco collaborativo con la seguente griglia come tabellone di gioco:} \\
\\ 
\texttt{```} \\
\texttt{ 1234567} \\
\texttt{╔═══╤═╤═╗} \\
\texttt{║◌◌I│W│◌║ 1} \\
\texttt{║◌◌◌└─┤◌║ 2} \\
\texttt{║◌◌◌TC├─╢ 3} \\
\texttt{╟─┐◌┌─┤◌║ 4} \\
\texttt{║◌│◌│◌│◌║ 5} \\
\texttt{╟─┼─┼─┼─╢ 6} \\
\texttt{║◌│◌│H│◌║ 7} \\
\texttt{╚═╧═╧═╧═╝} \\
\\ 
\texttt{```} \\
\\ 
\texttt{* Il bordo superiore mostra le coordinate x che aumentano verso destra, e il bordo destro le coordinate y che aumentano verso il basso.} \\
\texttt{* I seguenti oggetti sono posizionati casualmente sulla tua griglia: 'W', 'I', 'T', 'C', 'H'.} \\
\\ 
\\ 
\texttt{L'altro giocatore vede una variante del tabellone di gioco, dove gli oggetti sono posizionati in diverse posizioni casuali. Non puoi vedere il tabellone dell'altro giocatore e lui non può vedere il tuo.} \\
\\ 
\\ 
\texttt{**Obiettivo:**} \\
\\ 
\texttt{Entrambi i giocatori devono spostare gli oggetti sui rispettivi sfondi in modo che oggetti identici finiscano alle stesse coordinate. Devi comunicare con l'altro giocatore per concordare una configurazione obiettivo comune.} \\
\\ 
\\ 
\texttt{**Regole:**} \\
\\ 
\texttt{* In ogni turno, puoi inviare esattamente uno dei seguenti due comandi:} \\
\texttt{ 1. `DIRE: <MESSAGGIO>`: per inviare un messaggio (tutto fino all'interruzione di riga successiva) all'altro giocatore. Lo inoltrerò al tuo partner.} \\
\texttt{2. `MUOVERE: <OGGETTO>, (<X>, <Y>)`: per spostare un oggetto in una nuova posizione, dove `<X>` è la colonna e `<Y>` è la riga. Ti informerò se il tuo spostamento ha avuto successo o meno.} \\
\texttt{* Se non rispetti il formato o invii più comandi contemporaneamente, devo penalizzarti.} \\
\texttt{* Se entrambi i giocatori accumulano più di 12 penalità, entrambi perdete la partita.} \\
\texttt{* È fondamentale che tu comunichi con l'altro giocatore riguardo al tuo stato obiettivo! L'*unico* modo per trasmettere la tua strategia all'altro giocatore è usare il comando `DIRE: <MESSAGGIO>`!} \\
\\ 
\\ 
\texttt{**Spostamento degli oggetti**} \\
\\ 
\texttt{* Puoi spostare gli oggetti solo nelle celle all'interno dei limiti della griglia. La cella di destinazione deve essere vuota, cioè deve contenere solo il simbolo '◌'.} \\
\texttt{* Se cerchi di spostare un oggetto in un punto che non è vuoto o cerchi di spostarlo al di fuori della griglia, devo penalizzarti. Avrai un altro tentativo.} \\
\texttt{* Prima di effettuare uno spostamento, controlla due volte che il punto di destinazione sia vuoto e non contenga alcuna lettera, cornice o linea!} \\
\\ 
\\ 
\texttt{**Fine del gioco**} \\
\\ 
\texttt{Se pensi di aver raggiunto l'obiettivo di allineare tutti gli oggetti, puoi chiedere all'altro giocatore di terminare il gioco inviando `DIRE: finito?`. Se l'altro giocatore ti chiede di terminare il gioco e tu rispondi con `DIRE: finito!`, il gioco terminerà.} \\
\\ 
\texttt{Entrambi i giocatori vincono se il gioco termina entro 20 turni, dove un turno è definito come due giocatori che inviano ciascuno un comando valido.} \\
\\ 
\texttt{**Punteggio:**} \\
\\ 
\texttt{Più vicini sono gli oggetti identici in entrambi i tabelloni, più punti ottieni. Le penalità riducono i tuoi punti. Riuscirai a battere il record?} \\
\\ 
\\ 
\texttt{Invia un messaggio all'altro giocatore per iniziare la partita!} \\
            }
        }
    }
    & \\ \\

    \theutterance \stepcounter{utterance}  
    & \multicolumn{4}{p{0.6\linewidth}}{
        \cellcolor[rgb]{0.9,0.9,0.9}{
            \makecell[{{p{\linewidth}}}]{
                \texttt{\tiny{[P1$\rangle$GM]}}
                \texttt{DIRE: Ciao! Guardiamo le nostre posizioni iniziali. I miei oggetti sono:} \\
\texttt{W: (3,1)} \\
\texttt{I: (3,1) {[}Aspetta, controlliamo la griglia...{]}} \\
\\ 
\texttt{Rileggiamo la griglia con coordinate (x, y). X da 1 a 7, Y da 1 a 7.} \\
\\ 
\texttt{Riga 1 (y=1): `║◌◌I│W│◌║`} \\
\texttt{Colonne: 1=◌, 2=◌, 3=I, 4=│ (muro?), 5=W, 6=│?, 7=◌? No, la griglia è 7 colonne.} \\
\texttt{L'intestazione è `1234567`.} \\
\texttt{La riga 1: `║◌◌I│W│◌║`} \\
\texttt{Indici: 1:◌, 2:◌, 3:I, 4:│, 5:W, 6:│, 7:◌.} \\
\texttt{Ma ci sono muri verticali `│` che separano le celle? Sembra una griglia con muri interni.} \\
\texttt{Guardiamo la legenda: "La cella di destinazione deve essere vuota, cioè deve contenere solo il simbolo '◌'."} \\
\texttt{I muri (`│`, `─`, `└`, `├`, `┤`, `┐`, `┬`, `┴`, `┼`, `║`, `═`, `╤`, `╧`, `╔`, `╗`, `╚`, `╝`, `╟`, `╢`, `╠`, `╣`) sono ostacoli/fissi.} \\
\texttt{Le celle percorribili sono solo quelle con `◌` o le lettere (W, I, T, C, H).} \\
\\ 
\texttt{Analizziamo la mappa cella per cella (7x7).} \\
\texttt{Riga 1 (y=1): `║◌◌I│W│◌║` {-}> Celle x=1..7} \\
\texttt{x=1: ◌} \\
\texttt{x=2: ◌} \\
\texttt{x=3: I} \\
\texttt{x=4: │ (Muro)} \\
\texttt{x=5: W} \\
\texttt{x=6: │ (Muro)} \\
\texttt{x=7: ◌} \\
\\ 
\texttt{Riga 2 (y=2): `║◌◌◌└─┤◌║`} \\
\texttt{x=1: ◌} \\
\texttt{x=2: ◌} \\
\texttt{x=3: ◌} \\
\texttt{x=4: └ (Muro angolo)} \\
\texttt{x=5: ─ (Muro orizzontale)} \\
\texttt{x=6: ┤ (Muro T)} \\
\texttt{x=7: ◌} \\
\\ 
\texttt{Riga 3 (y=3): `║◌◌◌TC├─╢`} \\
\texttt{x=1: ◌} \\
\texttt{x=2: ◌} \\
\texttt{x=3: ◌} \\
\texttt{x=4: T} \\
\texttt{x=5: C} \\
\texttt{x=6: ├ (Muro T)} \\
\texttt{x=7: ─ (Muro) ? No, `╢` è bordo destro. `├─╢` : x=6 è `├`, x=7 è `─`? Ma il bordo è `╢`.} \\
\texttt{Aspetta, la stringa è `║◌◌◌TC├─╢`. Lunghezza caratteri: `║` (1), `◌`(1), `◌`(1), `◌`(1), `T`(1), `C`(1), `├`(1), `─`(1), `╢`(1). Totale 9 char per 7 colonne + 2 bordi.} \\
\texttt{Significa che ogni colonna è 1 char.} \\
\texttt{x=1: ◌} \\
\texttt{x=2: ◌} \\
\texttt{x=3: ◌} \\
\texttt{x=4: T} \\
\texttt{x=5: C} \\
\texttt{x=6: ├ (Muro)} \\
\texttt{x=7: ─ (Muro) {-}> Ma x=7 dovrebbe essere una cella. `╢` è l'angolo in basso a destra del bordo? No, `╢` è "box drawings vertical single and left double". È il bordo destro.} \\
\texttt{Quindi la riga 3 ha 7 celle: indices 1{-}7.} \\
\texttt{Contenuto: 1:◌, 2:◌, 3:◌, 4:T, 5:C, 6:├, 7:─? No, `├─╢` sono 3 caratteri per le ultime 2 colonne + bordo?} \\
\texttt{Guardiamo la riga 1: `║◌◌I│W│◌║` {-}> `║` + 7 char + `║`? `◌◌I│W│◌` sono 7 caratteri. Sì.} \\
\texttt{Riga 3: `◌◌◌TC├─` sono 7 caratteri? `◌` `◌` `◌` `T` `C` `├` `─`. Sì, 7 caratteri. Poi `╢` bordo.} \\
\texttt{Quindi:} \\
\texttt{Riga 3 (y=3):} \\
\texttt{x=1: ◌} \\
\texttt{x=2: ◌} \\
\texttt{x=3: ◌} \\
\texttt{x=4: T} \\
\texttt{x=5: C} \\
\texttt{x=6: ├ (Muro)} \\
\texttt{x=7: ─ (Muro orizzontale) {-}> Cella bloccata.} \\
\\ 
\texttt{Riga 4 (y=4): `╟─┐◌┌─┤◌║` {-}> Contenuto: `─┐◌┌─┤◌` (7 char)} \\
\texttt{x=1: ─ (Muro)} \\
\texttt{x=2: ┐ (Muro)} \\
\texttt{x=3: ◌} \\
\texttt{x=4: ┌ (Muro)} \\
\texttt{x=5: ─ (Muro)} \\
\texttt{x=6: ┤ (Muro)} \\
\texttt{x=7: ◌} \\
\\ 
\texttt{Riga 5 (y=5): `║◌│◌│◌│◌║` {-}> Contenuto: `◌│◌│◌│◌` (7 char)} \\
\texttt{x=1: ◌} \\
\texttt{x=2: │ (Muro)} \\
\texttt{x=3: ◌} \\
\texttt{x=4: │ (Muro)} \\
\texttt{x=5: ◌} \\
\texttt{x=6: │ (Muro)} \\
\texttt{x=7: ◌} \\
\\ 
\texttt{Riga 6 (y=6): `╟─┼─┼─┼─╢` {-}> Contenuto: `─┼─┼─┼─` (7 char) {-}> Tutti muri. Riga 6 completamente bloccata da muri orizzontali/incroci.} \\
\\ 
\texttt{Riga 7 (y=7): `║◌│◌│H│◌║` {-}> Contenuto: `◌│◌│H│◌` (7 char)} \\
\texttt{x=1: ◌} \\
\texttt{x=2: │ (Muro)} \\
\texttt{x=3: ◌} \\
\texttt{x=4: │ (Muro)} \\
\texttt{x=5: H} \\
\texttt{x=6: │ (Muro)} \\
\texttt{x=7: ◌} \\
\\ 
\texttt{Riassunto Oggetti MIEI:} \\
\texttt{I: (3, 1)} \\
\texttt{W: (5, 1)} \\
\texttt{T: (4, 3)} \\
\texttt{C: (5, 3)} \\
\texttt{H: (5, 7)} \\
\\ 
\texttt{Celle libere (◌) accessibili? Dobbiamo vedere la connettività. I muri bloccano il movimento? Le regole dicono: "spostare un oggetto in una nuova posizione... La cella di destinazione deve essere vuota". Non dice che il percorso deve essere libero, solo la destinazione. È un "teletrasporto" o uno scorrimento? "spostare un oggetto in una nuova posizione". Sembra teletrasporto in una cella vuota qualsiasi della griglia, ignorando i muri per il percorso, ma la cella destinazione deve essere `◌` (vuota). I muri occupano celle, quindi non si può mettere un oggetto su un muro.} \\
\\ 
\texttt{Quindi le coordinate valide per lo spostamento sono solo quelle che attualmente contengono `◌`.} \\
\texttt{Elenco celle `◌` (vuote):} \\
\texttt{(1,1), (2,1), (7,1)} \\
\texttt{(1,2), (2,2), (3,2), (7,2)} \\
\texttt{(1,3), (2,3), (3,3)} \\
\texttt{(3,4), (7,4)} \\
\texttt{(1,5), (3,5), (5,5), (7,5)} \\
\texttt{(1,7), (3,7), (7,7)} \\
\\ 
\texttt{Totale 18 celle vuote. 5 oggetti. Dobbiamo metterli d'accordo con l'altro giocatore.} \\
\\ 
\texttt{Strategia: Proporre una configurazione canonica. Ad esempio, angoli o una riga/colonna libera.} \\
\texttt{Guardiamo le celle vuote comuni? Non sappiamo la mappa dell'altro. Ma la struttura dei muri potrebbe essere} \\
            }
        }
    }
    & & \\ \\

    \theutterance \stepcounter{utterance}  
    & & & \multicolumn{2}{p{0.3\linewidth}}{
        \cellcolor[rgb]{0.9,0.9,0.9}{
            \makecell[{{p{\linewidth}}}]{
                \texttt{\tiny{[GM$|$GM]}}
                \texttt{DIRE: Ciao! Guardiamo le nostre posizioni iniziali. I miei oggetti sono:} \\
\texttt{W: (3,1)} \\
\texttt{I: (3,1) {[}Aspetta, controlliamo la griglia...{]}} \\
\\ 
\texttt{Rileggiamo la griglia con coordinate (x, y). X da 1 a 7, Y da 1 a 7.} \\
\\ 
\texttt{Riga 1 (y=1): `║◌◌I│W│◌║`} \\
\texttt{Colonne: 1=◌, 2=◌, 3=I, 4=│ (muro?), 5=W, 6=│?, 7=◌? No, la griglia è 7 colonne.} \\
\texttt{L'intestazione è `1234567`.} \\
\texttt{La riga 1: `║◌◌I│W│◌║`} \\
\texttt{Indici: 1:◌, 2:◌, 3:I, 4:│, 5:W, 6:│, 7:◌.} \\
\texttt{Ma ci sono muri verticali `│` che separano le celle? Sembra una griglia con muri interni.} \\
\texttt{Guardiamo la legenda: "La cella di destinazione deve essere vuota, cioè deve contenere solo il simbolo '◌'."} \\
\texttt{I muri (`│`, `─`, `└`, `├`, `┤`, `┐`, `┬`, `┴`, `┼`, `║`, `═`, `╤`, `╧`, `╔`, `╗`, `╚`, `╝`, `╟`, `╢`, `╠`, `╣`) sono ostacoli/fissi.} \\
\texttt{Le celle percorribili sono solo quelle con `◌` o le lettere (W, I, T, C, H).} \\
\\ 
\texttt{Analizziamo la mappa cella per cella (7x7).} \\
\texttt{Riga 1 (y=1): `║◌◌I│W│◌║` {-}> Celle x=1..7} \\
\texttt{x=1: ◌} \\
\texttt{x=2: ◌} \\
\texttt{x=3: I} \\
\texttt{x=4: │ (Muro)} \\
\texttt{x=5: W} \\
\texttt{x=6: │ (Muro)} \\
\texttt{x=7: ◌} \\
\\ 
\texttt{Riga 2 (y=2): `║◌◌◌└─┤◌║`} \\
\texttt{x=1: ◌} \\
\texttt{x=2: ◌} \\
\texttt{x=3: ◌} \\
\texttt{x=4: └ (Muro angolo)} \\
\texttt{x=5: ─ (Muro orizzontale)} \\
\texttt{x=6: ┤ (Muro T)} \\
\texttt{x=7: ◌} \\
\\ 
\texttt{Riga 3 (y=3): `║◌◌◌TC├─╢`} \\
\texttt{x=1: ◌} \\
\texttt{x=2: ◌} \\
\texttt{x=3: ◌} \\
\texttt{x=4: T} \\
\texttt{x=5: C} \\
\texttt{x=6: ├ (Muro T)} \\
\texttt{x=7: ─ (Muro) ? No, `╢` è bordo destro. `├─╢` : x=6 è `├`, x=7 è `─`? Ma il bordo è `╢`.} \\
\texttt{Aspetta, la stringa è `║◌◌◌TC├─╢`. Lunghezza caratteri: `║` (1), `◌`(1), `◌`(1), `◌`(1), `T`(1), `C`(1), `├`(1), `─`(1), `╢`(1). Totale 9 char per 7 colonne + 2 bordi.} \\
\texttt{Significa che ogni colonna è 1 char.} \\
\texttt{x=1: ◌} \\
\texttt{x=2: ◌} \\
\texttt{x=3: ◌} \\
\texttt{x=4: T} \\
\texttt{x=5: C} \\
\texttt{x=6: ├ (Muro)} \\
\texttt{x=7: ─ (Muro) {-}> Ma x=7 dovrebbe essere una cella. `╢` è l'angolo in basso a destra del bordo? No, `╢` è "box drawings vertical single and left double". È il bordo destro.} \\
\texttt{Quindi la riga 3 ha 7 celle: indices 1{-}7.} \\
\texttt{Contenuto: 1:◌, 2:◌, 3:◌, 4:T, 5:C, 6:├, 7:─? No, `├─╢` sono 3 caratteri per le ultime 2 colonne + bordo?} \\
\texttt{Guardiamo la riga 1: `║◌◌I│W│◌║` {-}> `║` + 7 char + `║`? `◌◌I│W│◌` sono 7 caratteri. Sì.} \\
\texttt{Riga 3: `◌◌◌TC├─` sono 7 caratteri? `◌` `◌` `◌` `T` `C` `├` `─`. Sì, 7 caratteri. Poi `╢` bordo.} \\
\texttt{Quindi:} \\
\texttt{Riga 3 (y=3):} \\
\texttt{x=1: ◌} \\
\texttt{x=2: ◌} \\
\texttt{x=3: ◌} \\
\texttt{x=4: T} \\
\texttt{x=5: C} \\
\texttt{x=6: ├ (Muro)} \\
\texttt{x=7: ─ (Muro orizzontale) {-}> Cella bloccata.} \\
\\ 
\texttt{Riga 4 (y=4): `╟─┐◌┌─┤◌║` {-}> Contenuto: `─┐◌┌─┤◌` (7 char)} \\
\texttt{x=1: ─ (Muro)} \\
\texttt{x=2: ┐ (Muro)} \\
\texttt{x=3: ◌} \\
\texttt{x=4: ┌ (Muro)} \\
\texttt{x=5: ─ (Muro)} \\
\texttt{x=6: ┤ (Muro)} \\
\texttt{x=7: ◌} \\
\\ 
\texttt{Riga 5 (y=5): `║◌│◌│◌│◌║` {-}> Contenuto: `◌│◌│◌│◌` (7 char)} \\
\texttt{x=1: ◌} \\
\texttt{x=2: │ (Muro)} \\
\texttt{x=3: ◌} \\
\texttt{x=4: │ (Muro)} \\
\texttt{x=5: ◌} \\
\texttt{x=6: │ (Muro)} \\
\texttt{x=7: ◌} \\
\\ 
\texttt{Riga 6 (y=6): `╟─┼─┼─┼─╢` {-}> Contenuto: `─┼─┼─┼─` (7 char) {-}> Tutti muri. Riga 6 completamente bloccata da muri orizzontali/incroci.} \\
\\ 
\texttt{Riga 7 (y=7): `║◌│◌│H│◌║` {-}> Contenuto: `◌│◌│H│◌` (7 char)} \\
\texttt{x=1: ◌} \\
\texttt{x=2: │ (Muro)} \\
\texttt{x=3: ◌} \\
\texttt{x=4: │ (Muro)} \\
\texttt{x=5: H} \\
\texttt{x=6: │ (Muro)} \\
\texttt{x=7: ◌} \\
\\ 
\texttt{Riassunto Oggetti MIEI:} \\
\texttt{I: (3, 1)} \\
\texttt{W: (5, 1)} \\
\texttt{T: (4, 3)} \\
\texttt{C: (5, 3)} \\
\texttt{H: (5, 7)} \\
\\ 
\texttt{Celle libere (◌) accessibili? Dobbiamo vedere la connettività. I muri bloccano il movimento? Le regole dicono: "spostare un oggetto in una nuova posizione... La cella di destinazione deve essere vuota". Non dice che il percorso deve essere libero, solo la destinazione. È un "teletrasporto" o uno scorrimento? "spostare un oggetto in una nuova posizione". Sembra teletrasporto in una cella vuota qualsiasi della griglia, ignorando i muri per il percorso, ma la cella destinazione deve essere `◌` (vuota). I muri occupano celle, quindi non si può mettere un oggetto su un muro.} \\
\\ 
\texttt{Quindi le coordinate valide per lo spostamento sono solo quelle che attualmente contengono `◌`.} \\
\texttt{Elenco celle `◌` (vuote):} \\
\texttt{(1,1), (2,1), (7,1)} \\
\texttt{(1,2), (2,2), (3,2), (7,2)} \\
\texttt{(1,3), (2,3), (3,3)} \\
\texttt{(3,4), (7,4)} \\
\texttt{(1,5), (3,5), (5,5), (7,5)} \\
\texttt{(1,7), (3,7), (7,7)} \\
\\ 
\texttt{Totale 18 celle vuote. 5 oggetti. Dobbiamo metterli d'accordo con l'altro giocatore.} \\
\\ 
\texttt{Strategia: Proporre una configurazione canonica. Ad esempio, angoli o una riga/colonna libera.} \\
\texttt{Guardiamo le celle vuote comuni? Non sappiamo la mappa dell'altro. Ma la struttura dei muri potrebbe essere} \\
            }
        }
    }
    & & \\ \\

    \theutterance \stepcounter{utterance}  
    & & & \multicolumn{2}{p{0.3\linewidth}}{
        \cellcolor[rgb]{0.9,0.9,0.9}{
            \makecell[{{p{\linewidth}}}]{
                \texttt{\tiny{[GM$|$GM]}}
                \texttt{Invalid format: tail is not empty: '} \\
\texttt{W: (3,1)} \\
\texttt{I: (3,1) {[}Aspetta, controlliamo la griglia...{]}} \\
\\ 
\texttt{Rileggiamo la griglia con coordinate (x, y). X da 1 a 7, Y da 1 a 7.} \\
\\ 
\texttt{Riga 1 (y=1): ║◌◌I│W│◌║} \\
\texttt{Colonne: 1=◌, 2=◌, 3=I, 4=│ (muro?), 5=W, 6=│?, 7=◌? No, la griglia è 7 colonne.} \\
\texttt{L'intestazione è 1234567.} \\
\texttt{La riga 1: ║◌◌I│W│◌║} \\
\texttt{Indici: 1:◌, 2:◌, 3:I, 4:│, 5:W, 6:│, 7:◌.} \\
\texttt{Ma ci sono muri verticali │ che separano le celle? Sembra una griglia con muri interni.} \\
\texttt{Guardiamo la legenda: "La cella di destinazione deve essere vuota, cioè deve contenere solo il simbolo '◌'."} \\
\texttt{I muri (│, ─, └, ├, ┤, ┐, ┬, ┴, ┼, ║, ═, ╤, ╧, ╔, ╗, ╚, ╝, ╟, ╢, ╠, ╣) sono ostacoli/fissi.} \\
\texttt{Le celle percorribili sono solo quelle con ◌ o le lettere (W, I, T, C, H).} \\
\\ 
\texttt{Analizziamo la mappa cella per cella (7x7).} \\
\texttt{Riga 1 (y=1): ║◌◌I│W│◌║ {-}> Celle x=1..7} \\
\texttt{x=1: ◌} \\
\texttt{x=2: ◌} \\
\texttt{x=3: I} \\
\texttt{x=4: │ (Muro)} \\
\texttt{x=5: W} \\
\texttt{x=6: │ (Muro)} \\
\texttt{x=7: ◌} \\
\\ 
\texttt{Riga 2 (y=2): ║◌◌◌└─┤◌║} \\
\texttt{x=1: ◌} \\
\texttt{x=2: ◌} \\
\texttt{x=3: ◌} \\
\texttt{x=4: └ (Muro angolo)} \\
\texttt{x=5: ─ (Muro orizzontale)} \\
\texttt{x=6: ┤ (Muro T)} \\
\texttt{x=7: ◌} \\
\\ 
\texttt{Riga 3 (y=3): ║◌◌◌TC├─╢} \\
\texttt{x=1: ◌} \\
\texttt{x=2: ◌} \\
\texttt{x=3: ◌} \\
\texttt{x=4: T} \\
\texttt{x=5: C} \\
\texttt{x=6: ├ (Muro T)} \\
\texttt{x=7: ─ (Muro) ? No, ╢ è bordo destro. ├─╢ : x=6 è ├, x=7 è ─? Ma il bordo è ╢.} \\
\texttt{Aspetta, la stringa è ║◌◌◌TC├─╢. Lunghezza caratteri: ║ (1), ◌(1), ◌(1), ◌(1), T(1), C(1), ├(1), ─(1), ╢(1). Totale 9 char per 7 colonne + 2 bordi.} \\
\texttt{Significa che ogni colonna è 1 char.} \\
\texttt{x=1: ◌} \\
\texttt{x=2: ◌} \\
\texttt{x=3: ◌} \\
\texttt{x=4: T} \\
\texttt{x=5: C} \\
\texttt{x=6: ├ (Muro)} \\
\texttt{x=7: ─ (Muro) {-}> Ma x=7 dovrebbe essere una cella. ╢ è l'angolo in basso a destra del bordo? No, ╢ è "box drawings vertical single and left double". È il bordo destro.} \\
\texttt{Quindi la riga 3 ha 7 celle: indices 1{-}7.} \\
\texttt{Contenuto: 1:◌, 2:◌, 3:◌, 4:T, 5:C, 6:├, 7:─? No, ├─╢ sono 3 caratteri per le ultime 2 colonne + bordo?} \\
\texttt{Guardiamo la riga 1: ║◌◌I│W│◌║ {-}> ║ + 7 char + ║? ◌◌I│W│◌ sono 7 caratteri. Sì.} \\
\texttt{Riga 3: ◌◌◌TC├─ sono 7 caratteri? ◌ ◌ ◌ T C ├ ─. Sì, 7 caratteri. Poi ╢ bordo.} \\
\texttt{Quindi:} \\
\texttt{Riga 3 (y=3):} \\
\texttt{x=1: ◌} \\
\texttt{x=2: ◌} \\
\texttt{x=3: ◌} \\
\texttt{x=4: T} \\
\texttt{x=5: C} \\
\texttt{x=6: ├ (Muro)} \\
\texttt{x=7: ─ (Muro orizzontale) {-}> Cella bloccata.} \\
\\ 
\texttt{Riga 4 (y=4): ╟─┐◌┌─┤◌║ {-}> Contenuto: ─┐◌┌─┤◌ (7 char)} \\
\texttt{x=1: ─ (Muro)} \\
\texttt{x=2: ┐ (Muro)} \\
\texttt{x=3: ◌} \\
\texttt{x=4: ┌ (Muro)} \\
\texttt{x=5: ─ (Muro)} \\
\texttt{x=6: ┤ (Muro)} \\
\texttt{x=7: ◌} \\
\\ 
\texttt{Riga 5 (y=5): ║◌│◌│◌│◌║ {-}> Contenuto: ◌│◌│◌│◌ (7 char)} \\
\texttt{x=1: ◌} \\
\texttt{x=2: │ (Muro)} \\
\texttt{x=3: ◌} \\
\texttt{x=4: │ (Muro)} \\
\texttt{x=5: ◌} \\
\texttt{x=6: │ (Muro)} \\
\texttt{x=7: ◌} \\
\\ 
\texttt{Riga 6 (y=6): ╟─┼─┼─┼─╢ {-}> Contenuto: ─┼─┼─┼─ (7 char) {-}> Tutti muri. Riga 6 completamente bloccata da muri orizzontali/incroci.} \\
\\ 
\texttt{Riga 7 (y=7): ║◌│◌│H│◌║ {-}> Contenuto: ◌│◌│H│◌ (7 char)} \\
\texttt{x=1: ◌} \\
\texttt{x=2: │ (Muro)} \\
\texttt{x=3: ◌} \\
\texttt{x=4: │ (Muro)} \\
\texttt{x=5: H} \\
\texttt{x=6: │ (Muro)} \\
\texttt{x=7: ◌} \\
\\ 
\texttt{Riassunto Oggetti MIEI:} \\
\texttt{I: (3, 1)} \\
\texttt{W: (5, 1)} \\
\texttt{T: (4, 3)} \\
\texttt{C: (5, 3)} \\
\texttt{H: (5, 7)} \\
\\ 
\texttt{Celle libere (◌) accessibili? Dobbiamo vedere la connettività. I muri bloccano il movimento? Le regole dicono: "spostare un oggetto in una nuova posizione... La cella di destinazione deve essere vuota". Non dice che il percorso deve essere libero, solo la destinazione. È un "teletrasporto" o uno scorrimento? "spostare un oggetto in una nuova posizione". Sembra teletrasporto in una cella vuota qualsiasi della griglia, ignorando i muri per il percorso, ma la cella destinazione deve essere ◌ (vuota). I muri occupano celle, quindi non si può mettere un oggetto su un muro.} \\
\\ 
\texttt{Quindi le coordinate valide per lo spostamento sono solo quelle che attualmente contengono ◌.} \\
\texttt{Elenco celle ◌ (vuote):} \\
\texttt{(1,1), (2,1), (7,1)} \\
\texttt{(1,2), (2,2), (3,2), (7,2)} \\
\texttt{(1,3), (2,3), (3,3)} \\
\texttt{(3,4), (7,4)} \\
\texttt{(1,5), (3,5), (5,5), (7,5)} \\
\texttt{(1,7), (3,7), (7,7)} \\
\\ 
\texttt{Totale 18 celle vuote. 5 oggetti. Dobbiamo metterli d'accordo con l'altro giocatore.} \\
\\ 
\texttt{Strategia: Proporre una configurazione canonica. Ad esempio, angoli o una riga/colonna libera.} \\
\texttt{Guardiamo le celle vuote comuni? Non sappiamo la mappa dell'altro. Ma la struttura dei muri potrebbe essere} \\
            }
        }
    }
    & & \\ \\

    \theutterance \stepcounter{utterance}  
    & & \multicolumn{4}{p{0.6\linewidth}}{
        \cellcolor[rgb]{0.9,0.9,0.9}{
            \makecell[{{p{\linewidth}}}]{
                \texttt{\tiny{[P1$\langle$GM]}}
                \texttt{Penalità: Il tuo messaggio non deve contenere nulla dopo il comando!} \\
\\ 
\texttt{Assicurati che la tua risposta contenga solo `DIRE: <MESSAGGIO>` o `MUOVERE: <OGGETTO>, (<X>, <Y>)`, e nient'altro!} \\
\\ 
\texttt{Avete accumulato collettivamente 1 di 12 penalità.} \\
\texttt{Per favore riprova!} \\
            }
        }
    }
    & \\ \\

    \theutterance \stepcounter{utterance}  
    & \multicolumn{4}{p{0.6\linewidth}}{
        \cellcolor[rgb]{0.9,0.9,0.9}{
            \makecell[{{p{\linewidth}}}]{
                \texttt{\tiny{[P1$\rangle$GM]}}
                \texttt{DIRE: Ciao! Ecco le mie posizioni iniziali: I(3,1), W(5,1), T(4,3), C(5,3), H(5,7). Celle vuote disponibili (solo ◌): (1,1),(2,1),(7,1), (1,2),(2,2),(3,2),(7,2), (1,3),(2,3),(3,3), (3,4),(7,4), (1,5),(3,5),(5,5),(7,5), (1,7),(3,7),(7,7). Quali sono le tue posizioni? Proponiamo un obiettivo comune.} \\
            }
        }
    }
    & & \\ \\

    \theutterance \stepcounter{utterance}  
    & & & \multicolumn{2}{p{0.3\linewidth}}{
        \cellcolor[rgb]{0.9,0.9,0.9}{
            \makecell[{{p{\linewidth}}}]{
                \texttt{\tiny{[GM$|$GM]}}
                \texttt{DIRE: Ciao! Ecco le mie posizioni iniziali: I(3,1), W(5,1), T(4,3), C(5,3), H(5,7). Celle vuote disponibili (solo ◌): (1,1),(2,1),(7,1), (1,2),(2,2),(3,2),(7,2), (1,3),(2,3),(3,3), (3,4),(7,4), (1,5),(3,5),(5,5),(7,5), (1,7),(3,7),(7,7). Quali sono le tue posizioni? Proponiamo un obiettivo comune.} \\
            }
        }
    }
    & & \\ \\

    \theutterance \stepcounter{utterance}  
    & & \multicolumn{4}{p{0.6\linewidth}}{
        \cellcolor[rgb]{0.9,0.9,0.9}{
            \makecell[{{p{\linewidth}}}]{
                \texttt{\tiny{[GM$\rangle$P2]}}
                \texttt{Sono il tuo maestro di gioco e stai giocando un gioco collaborativo con la seguente griglia come tabellone di gioco:} \\
\\ 
\texttt{```} \\
\texttt{ 1234567} \\
\texttt{╔═══╤═╤═╗} \\
\texttt{║◌◌◌│◌│◌║ 1} \\
\texttt{║◌W◌└─┤C║ 2} \\
\texttt{║T◌◌◌◌├─╢ 3} \\
\texttt{╟─┐◌┌─┤◌║ 4} \\
\texttt{║◌│◌│H│I║ 5} \\
\texttt{╟─┼─┼─┼─╢ 6} \\
\texttt{║◌│◌│◌│◌║ 7} \\
\texttt{╚═╧═╧═╧═╝} \\
\\ 
\texttt{```} \\
\\ 
\texttt{* Il bordo superiore mostra le coordinate x che aumentano verso destra, e il bordo destro le coordinate y che aumentano verso il basso.} \\
\texttt{* I seguenti oggetti sono posizionati casualmente sulla tua griglia: 'W', 'I', 'T', 'C', 'H'.} \\
\\ 
\\ 
\texttt{L'altro giocatore vede una variante del tabellone di gioco, dove gli oggetti sono posizionati in diverse posizioni casuali. Non puoi vedere il tabellone dell'altro giocatore e lui non può vedere il tuo.} \\
\\ 
\\ 
\texttt{**Obiettivo:**} \\
\\ 
\texttt{Entrambi i giocatori devono spostare gli oggetti sui rispettivi sfondi in modo che oggetti identici finiscano alle stesse coordinate. Devi comunicare con l'altro giocatore per concordare una configurazione obiettivo comune.} \\
\\ 
\\ 
\texttt{**Regole:**} \\
\\ 
\texttt{* In ogni turno, puoi inviare esattamente uno dei seguenti due comandi:} \\
\texttt{ 1. `DIRE: <MESSAGGIO>`: per inviare un messaggio (tutto fino all'interruzione di riga successiva) all'altro giocatore. Lo inoltrerò al tuo partner.} \\
\texttt{2. `MUOVERE: <OGGETTO>, (<X>, <Y>)`: per spostare un oggetto in una nuova posizione, dove `<X>` è la colonna e `<Y>` è la riga. Ti informerò se il tuo spostamento ha avuto successo o meno.} \\
\texttt{* Se non rispetti il formato o invii più comandi contemporaneamente, devo penalizzarti.} \\
\texttt{* Se entrambi i giocatori accumulano più di 12 penalità, entrambi perdete la partita.} \\
\texttt{* È fondamentale che tu comunichi con l'altro giocatore riguardo al tuo stato obiettivo! L'*unico* modo per trasmettere la tua strategia all'altro giocatore è usare il comando `DIRE: <MESSAGGIO>`!} \\
\\ 
\\ 
\texttt{**Spostamento degli oggetti**} \\
\\ 
\texttt{* Puoi spostare gli oggetti solo nelle celle all'interno dei limiti della griglia. La cella di destinazione deve essere vuota, cioè deve contenere solo il simbolo '◌'.} \\
\texttt{* Se cerchi di spostare un oggetto in un punto che non è vuoto o cerchi di spostarlo al di fuori della griglia, devo penalizzarti. Avrai un altro tentativo.} \\
\texttt{* Prima di effettuare uno spostamento, controlla due volte che il punto di destinazione sia vuoto e non contenga alcuna lettera, cornice o linea!} \\
\\ 
\\ 
\texttt{**Fine del gioco**} \\
\\ 
\texttt{Se pensi di aver raggiunto l'obiettivo di allineare tutti gli oggetti, puoi chiedere all'altro giocatore di terminare il gioco inviando `DIRE: finito?`. Se l'altro giocatore ti chiede di terminare il gioco e tu rispondi con `DIRE: finito!`, il gioco terminerà.} \\
\\ 
\texttt{Entrambi i giocatori vincono se il gioco termina entro 20 turni, dove un turno è definito come due giocatori che inviano ciascuno un comando valido.} \\
\\ 
\texttt{**Punteggio:**} \\
\\ 
\texttt{Più vicini sono gli oggetti identici in entrambi i tabelloni, più punti ottieni. Le penalità riducono i tuoi punti. Riuscirai a battere il record?} \\
\\ 
\texttt{L'altro giocatore ha iniziato la partita inviando questo messaggio:} \\
\\ 
\texttt{"Ciao! Ecco le mie posizioni iniziali: I(3,1), W(5,1), T(4,3), C(5,3), H(5,7). Celle vuote disponibili (solo ◌): (1,1),(2,1),(7,1), (1,2),(2,2),(3,2),(7,2), (1,3),(2,3),(3,3), (3,4),(7,4), (1,5),(3,5),(5,5),(7,5), (1,7),(3,7),(7,7). Quali sono le tue posizioni? Proponiamo un obiettivo comune."} \\
\\ 
\texttt{Qual è il tuo primo comando?} \\
            }
        }
    }
    & \\ \\

    \theutterance \stepcounter{utterance}  
    & & & \multicolumn{4}{p{0.6\linewidth}}{
        \cellcolor[rgb]{0.9,0.9,0.9}{
            \makecell[{{p{\linewidth}}}]{
                \texttt{\tiny{[GM$\langle$P2]}}
                \texttt{DIRE: Ciao! Le mie posizioni iniziali sono: W(2,2), T(1,3), C(7,2), H(5,5), I(7,5). Celle vuote disponibili (solo ◌): (1,1),(3,1),(4,1),(5,1),(6,1), (1,2),(3,2),(4,2),(5,2),(6,2), (2,3),(3,3),(4,3),(5,3),(6,3), (3,4),(4,4),(6,4),(7,4), (1,5),(2,5),(3,5),(4,5),(6,5), (1,6),(2,6),(3,6),(4,6),(5,6),(6,6),(7,6), (1,7),(2,7),(3,7),(4,7),(5,7),(6,7).} \\
\\ 
\texttt{Poiché tu hai molte celle libere in alto (y=1,2) e io in basso (y=6,7), propongo di incontrarci a metà strada sulle righe 3 e 4 che sono libere per entrambi. Obiettivo proposto:} \\
\texttt{W {-}> (3,3)} \\
\texttt{T {-}> (4,3)} \\
\texttt{C {-}> (5,3)} \\
\texttt{H {-}> (5,4)} \\
\texttt{I {-}> (6,3)} \\
\\ 
\texttt{Tutte queste celle sono ◌ nella mia griglia. Confermi che sono libere anche nella tua? Se sì, inizio a muovere.} \\
            }
        }
    }
     \\ \\

    \theutterance \stepcounter{utterance}  
    & & & \multicolumn{2}{p{0.3\linewidth}}{
        \cellcolor[rgb]{0.9,0.9,0.9}{
            \makecell[{{p{\linewidth}}}]{
                \texttt{\tiny{[GM$|$GM]}}
                \texttt{DIRE: Ciao! Le mie posizioni iniziali sono: W(2,2), T(1,3), C(7,2), H(5,5), I(7,5). Celle vuote disponibili (solo ◌): (1,1),(3,1),(4,1),(5,1),(6,1), (1,2),(3,2),(4,2),(5,2),(6,2), (2,3),(3,3),(4,3),(5,3),(6,3), (3,4),(4,4),(6,4),(7,4), (1,5),(2,5),(3,5),(4,5),(6,5), (1,6),(2,6),(3,6),(4,6),(5,6),(6,6),(7,6), (1,7),(2,7),(3,7),(4,7),(5,7),(6,7).} \\
\\ 
\texttt{Poiché tu hai molte celle libere in alto (y=1,2) e io in basso (y=6,7), propongo di incontrarci a metà strada sulle righe 3 e 4 che sono libere per entrambi. Obiettivo proposto:} \\
\texttt{W {-}> (3,3)} \\
\texttt{T {-}> (4,3)} \\
\texttt{C {-}> (5,3)} \\
\texttt{H {-}> (5,4)} \\
\texttt{I {-}> (6,3)} \\
\\ 
\texttt{Tutte queste celle sono ◌ nella mia griglia. Confermi che sono libere anche nella tua? Se sì, inizio a muovere.} \\
            }
        }
    }
    & & \\ \\

    \theutterance \stepcounter{utterance}  
    & & & \multicolumn{2}{p{0.3\linewidth}}{
        \cellcolor[rgb]{0.9,0.9,0.9}{
            \makecell[{{p{\linewidth}}}]{
                \texttt{\tiny{[GM$|$GM]}}
                \texttt{Invalid format: tail is not empty: '} \\
\\ 
\texttt{Poiché tu hai molte celle libere in alto (y=1,2) e io in basso (y=6,7), propongo di incontrarci a metà strada sulle righe 3 e 4 che sono libere per entrambi. Obiettivo proposto:} \\
\texttt{W {-}> (3,3)} \\
\texttt{T {-}> (4,3)} \\
\texttt{C {-}> (5,3)} \\
\texttt{H {-}> (5,4)} \\
\texttt{I {-}> (6,3)} \\
\\ 
\texttt{Tutte queste celle sono ◌ nella mia griglia. Confermi che sono libere anche nella tua? Se sì, inizio a muovere.} \\
            }
        }
    }
    & & \\ \\

    \theutterance \stepcounter{utterance}  
    & & \multicolumn{4}{p{0.6\linewidth}}{
        \cellcolor[rgb]{0.9,0.9,0.9}{
            \makecell[{{p{\linewidth}}}]{
                \texttt{\tiny{[GM$\rangle$P2]}}
                \texttt{Penalità: Il tuo messaggio non deve contenere nulla dopo il comando!} \\
\\ 
\texttt{Assicurati che la tua risposta contenga solo `DIRE: <MESSAGGIO>` o `MUOVERE: <OGGETTO>, (<X>, <Y>)`, e nient'altro!} \\
\\ 
\texttt{Avete accumulato collettivamente 2 di 12 penalità.} \\
\texttt{Per favore riprova!} \\
            }
        }
    }
    & \\ \\

    \theutterance \stepcounter{utterance}  
    & & & \multicolumn{4}{p{0.6\linewidth}}{
        \cellcolor[rgb]{0.9,0.9,0.9}{
            \makecell[{{p{\linewidth}}}]{
                \texttt{\tiny{[GM$\langle$P2]}}
                \texttt{DIRE: Ciao! Le mie posizioni iniziali sono: W(2,2), T(1,3), C(7,2), H(5,5), I(7,5). Celle vuote disponibili: (1,1),(3,1),(4,1),(5,1),(6,1), (1,2),(3,2),(4,2),(5,2),(6,2), (2,3),(3,3),(4,3),(5,3),(6,3), (3,4),(4,4),(6,4),(7,4), (1,5),(2,5),(3,5),(4,5),(6,5), (1,6),(2,6),(3,6),(4,6),(5,6),(6,6),(7,6), (1,7),(2,7),(3,7),(4,7),(5,7),(6,7). Propongo obiettivo comune su righe 3{-}4: W(3,3), T(4,3), C(5,3), H(5,4), I(6,3). Sono libere anche nella tua griglia?} \\
            }
        }
    }
     \\ \\

    \theutterance \stepcounter{utterance}  
    & & & \multicolumn{2}{p{0.3\linewidth}}{
        \cellcolor[rgb]{0.9,0.9,0.9}{
            \makecell[{{p{\linewidth}}}]{
                \texttt{\tiny{[GM$|$GM]}}
                \texttt{DIRE: Ciao! Le mie posizioni iniziali sono: W(2,2), T(1,3), C(7,2), H(5,5), I(7,5). Celle vuote disponibili: (1,1),(3,1),(4,1),(5,1),(6,1), (1,2),(3,2),(4,2),(5,2),(6,2), (2,3),(3,3),(4,3),(5,3),(6,3), (3,4),(4,4),(6,4),(7,4), (1,5),(2,5),(3,5),(4,5),(6,5), (1,6),(2,6),(3,6),(4,6),(5,6),(6,6),(7,6), (1,7),(2,7),(3,7),(4,7),(5,7),(6,7). Propongo obiettivo comune su righe 3{-}4: W(3,3), T(4,3), C(5,3), H(5,4), I(6,3). Sono libere anche nella tua griglia?} \\
            }
        }
    }
    & & \\ \\

    \theutterance \stepcounter{utterance}  
    & & \multicolumn{4}{p{0.6\linewidth}}{
        \cellcolor[rgb]{0.9,0.9,0.9}{
            \makecell[{{p{\linewidth}}}]{
                \texttt{\tiny{[P1$\langle$GM]}}
                \texttt{Il tuo messaggio è stato trasmesso all'altro giocatore.} \\
\\ 
\\ 
\texttt{Stai giocando il round 1 di un massimo di 20.} \\
\texttt{Avete accumulato collettivamente 2 di 12 penalità.} \\
\texttt{L'altro giocatore ha inviato questo messaggio:} \\
\\ 
\\ 
\texttt{"Ciao! Le mie posizioni iniziali sono: W(2,2), T(1,3), C(7,2), H(5,5), I(7,5). Celle vuote disponibili: (1,1),(3,1),(4,1),(5,1),(6,1), (1,2),(3,2),(4,2),(5,2),(6,2), (2,3),(3,3),(4,3),(5,3),(6,3), (3,4),(4,4),(6,4),(7,4), (1,5),(2,5),(3,5),(4,5),(6,5), (1,6),(2,6),(3,6),(4,6),(5,6),(6,6),(7,6), (1,7),(2,7),(3,7),(4,7),(5,7),(6,7). Propongo obiettivo comune su righe 3{-}4: W(3,3), T(4,3), C(5,3), H(5,4), I(6,3). Sono libere anche nella tua griglia?"} \\
\\ 
\\ 
\texttt{Qual è il tuo prossimo comando?} \\
            }
        }
    }
    & \\ \\

    \theutterance \stepcounter{utterance}  
    & \multicolumn{4}{p{0.6\linewidth}}{
        \cellcolor[rgb]{0.9,0.9,0.9}{
            \makecell[{{p{\linewidth}}}]{
                \texttt{\tiny{[P1$\rangle$GM]}}
                \texttt{DIRE: Nella mia griglia, le celle (3,3), (4,3), (5,3), (5,4), (6,3) NON sono tutte libere: (3,3) e' libera, (4,3) ha T, (5,3) ha C, (6,3) e' muro (├), (5,4) e' muro (─). La mia riga 3 e' bloccata a destra. La mia riga 4 ha muri in (1,4),(2,4),(4,4),(5,4),(6,4). Riga 5 ha muri in colonne pari. Le mie celle libere "centrali" sono poche: (3,2), (3,4), (3,5). Propongo di usare la COLONNA 3 (x=3) che e' quasi tutta libera da y=1 a y=7 (eccetto y=1 occupato da I, y=3 libero, y=4 libero, y=5 libero, y=7 libero). Coordinate libere col 3: (3,2), (3,3), (3,4), (3,5), (3,7). Sono 5 celle perfette per 5 oggetti. Nella tua griglia, la colonna 3 (x=3) e' libera per y=1,2,3,4,5,6,7? Le tue celle libere includono (3,1),(3,2),(3,3),(3,4),(3,5),(3,6),(3,7). SI! Allineiamo tutto su x=3. Propongo: W(3,2), I(3,3), T(3,4), C(3,5), H(3,7). Confermi?} \\
            }
        }
    }
    & & \\ \\

    \theutterance \stepcounter{utterance}  
    & & & \multicolumn{2}{p{0.3\linewidth}}{
        \cellcolor[rgb]{0.9,0.9,0.9}{
            \makecell[{{p{\linewidth}}}]{
                \texttt{\tiny{[GM$|$GM]}}
                \texttt{DIRE: Nella mia griglia, le celle (3,3), (4,3), (5,3), (5,4), (6,3) NON sono tutte libere: (3,3) e' libera, (4,3) ha T, (5,3) ha C, (6,3) e' muro (├), (5,4) e' muro (─). La mia riga 3 e' bloccata a destra. La mia riga 4 ha muri in (1,4),(2,4),(4,4),(5,4),(6,4). Riga 5 ha muri in colonne pari. Le mie celle libere "centrali" sono poche: (3,2), (3,4), (3,5). Propongo di usare la COLONNA 3 (x=3) che e' quasi tutta libera da y=1 a y=7 (eccetto y=1 occupato da I, y=3 libero, y=4 libero, y=5 libero, y=7 libero). Coordinate libere col 3: (3,2), (3,3), (3,4), (3,5), (3,7). Sono 5 celle perfette per 5 oggetti. Nella tua griglia, la colonna 3 (x=3) e' libera per y=1,2,3,4,5,6,7? Le tue celle libere includono (3,1),(3,2),(3,3),(3,4),(3,5),(3,6),(3,7). SI! Allineiamo tutto su x=3. Propongo: W(3,2), I(3,3), T(3,4), C(3,5), H(3,7). Confermi?} \\
            }
        }
    }
    & & \\ \\

    \theutterance \stepcounter{utterance}  
    & & \multicolumn{4}{p{0.6\linewidth}}{
        \cellcolor[rgb]{0.9,0.9,0.9}{
            \makecell[{{p{\linewidth}}}]{
                \texttt{\tiny{[GM$\rangle$P2]}}
                \texttt{Il tuo messaggio è stato trasmesso all'altro giocatore.} \\
\\ 
\\ 
\texttt{Stai giocando il round 2 di un massimo di 20.} \\
\texttt{Avete accumulato collettivamente 2 di 12 penalità.} \\
\texttt{L'altro giocatore ha inviato questo messaggio:} \\
\\ 
\\ 
\texttt{"Nella mia griglia, le celle (3,3), (4,3), (5,3), (5,4), (6,3) NON sono tutte libere: (3,3) e' libera, (4,3) ha T, (5,3) ha C, (6,3) e' muro (├), (5,4) e' muro (─). La mia riga 3 e' bloccata a destra. La mia riga 4 ha muri in (1,4),(2,4),(4,4),(5,4),(6,4). Riga 5 ha muri in colonne pari. Le mie celle libere "centrali" sono poche: (3,2), (3,4), (3,5). Propongo di usare la COLONNA 3 (x=3) che e' quasi tutta libera da y=1 a y=7 (eccetto y=1 occupato da I, y=3 libero, y=4 libero, y=5 libero, y=7 libero). Coordinate libere col 3: (3,2), (3,3), (3,4), (3,5), (3,7). Sono 5 celle perfette per 5 oggetti. Nella tua griglia, la colonna 3 (x=3) e' libera per y=1,2,3,4,5,6,7? Le tue celle libere includono (3,1),(3,2),(3,3),(3,4),(3,5),(3,6),(3,7). SI! Allineiamo tutto su x=3. Propongo: W(3,2), I(3,3), T(3,4), C(3,5), H(3,7). Confermi?"} \\
\\ 
\\ 
\texttt{Qual è il tuo prossimo comando?} \\
            }
        }
    }
    & \\ \\

    \theutterance \stepcounter{utterance}  
    & & & \multicolumn{4}{p{0.6\linewidth}}{
        \cellcolor[rgb]{0.9,0.9,0.9}{
            \makecell[{{p{\linewidth}}}]{
                \texttt{\tiny{[GM$\langle$P2]}}
                \texttt{DIRE: Confermo! Nella mia griglia la colonna 3 (x=3) è completamente libera per y=1,2,3,4,5,6,7. Accetto l'obiettivo comune: W(3,2), I(3,3), T(3,4), C(3,5), H(3,7). Inizio a muovere i miei oggetti.} \\
            }
        }
    }
     \\ \\

    \theutterance \stepcounter{utterance}  
    & & & \multicolumn{2}{p{0.3\linewidth}}{
        \cellcolor[rgb]{0.9,0.9,0.9}{
            \makecell[{{p{\linewidth}}}]{
                \texttt{\tiny{[GM$|$GM]}}
                \texttt{DIRE: Confermo! Nella mia griglia la colonna 3 (x=3) è completamente libera per y=1,2,3,4,5,6,7. Accetto l'obiettivo comune: W(3,2), I(3,3), T(3,4), C(3,5), H(3,7). Inizio a muovere i miei oggetti.} \\
            }
        }
    }
    & & \\ \\

    \theutterance \stepcounter{utterance}  
    & & \multicolumn{4}{p{0.6\linewidth}}{
        \cellcolor[rgb]{0.9,0.9,0.9}{
            \makecell[{{p{\linewidth}}}]{
                \texttt{\tiny{[P1$\langle$GM]}}
                \texttt{Il tuo messaggio è stato trasmesso all'altro giocatore.} \\
\\ 
\\ 
\texttt{Stai giocando il round 2 di un massimo di 20.} \\
\texttt{Avete accumulato collettivamente 2 di 12 penalità.} \\
\texttt{L'altro giocatore ha inviato questo messaggio:} \\
\\ 
\\ 
\texttt{"Confermo! Nella mia griglia la colonna 3 (x=3) è completamente libera per y=1,2,3,4,5,6,7. Accetto l'obiettivo comune: W(3,2), I(3,3), T(3,4), C(3,5), H(3,7). Inizio a muovere i miei oggetti."} \\
\\ 
\\ 
\texttt{Qual è il tuo prossimo comando?} \\
            }
        }
    }
    & \\ \\

    \theutterance \stepcounter{utterance}  
    & \multicolumn{4}{p{0.6\linewidth}}{
        \cellcolor[rgb]{0.9,0.9,0.9}{
            \makecell[{{p{\linewidth}}}]{
                \texttt{\tiny{[P1$\rangle$GM]}}
                \texttt{MUOVERE: W, (3, 2)} \\
            }
        }
    }
    & & \\ \\

    \theutterance \stepcounter{utterance}  
    & & & \multicolumn{2}{p{0.3\linewidth}}{
        \cellcolor[rgb]{0.9,0.9,0.9}{
            \makecell[{{p{\linewidth}}}]{
                \texttt{\tiny{[GM$|$GM]}}
                \texttt{MUOVERE: W, (3, 2)} \\
            }
        }
    }
    & & \\ \\

    \theutterance \stepcounter{utterance}  
    & & & \multicolumn{2}{p{0.3\linewidth}}{
        \cellcolor[rgb]{0.9,0.9,0.9}{
            \makecell[{{p{\linewidth}}}]{
                \texttt{\tiny{[GM$|$GM]}}
                \texttt{W moved to (3, 2)} \\
            }
        }
    }
    & & \\ \\

    \theutterance \stepcounter{utterance}  
    & & \multicolumn{4}{p{0.6\linewidth}}{
        \cellcolor[rgb]{0.9,0.9,0.9}{
            \makecell[{{p{\linewidth}}}]{
                \texttt{\tiny{[GM$\rangle$P2]}}
                \texttt{Il tuo messaggio è stato trasmesso all'altro giocatore.} \\
\\ 
\\ 
\texttt{Stai giocando il round 3 di un massimo di 20.} \\
\texttt{Avete accumulato collettivamente 2 di 12 penalità.} \\
\texttt{L'altro giocatore ha spostato un oggetto sulla sua griglia.} \\
\\ 
\\ 
\texttt{Qual è il tuo prossimo comando?} \\
            }
        }
    }
    & \\ \\

    \theutterance \stepcounter{utterance}  
    & & & \multicolumn{4}{p{0.6\linewidth}}{
        \cellcolor[rgb]{0.9,0.9,0.9}{
            \makecell[{{p{\linewidth}}}]{
                \texttt{\tiny{[GM$\langle$P2]}}
                \texttt{MUOVERE: W, (3, 2)} \\
            }
        }
    }
     \\ \\

    \theutterance \stepcounter{utterance}  
    & & & \multicolumn{2}{p{0.3\linewidth}}{
        \cellcolor[rgb]{0.9,0.9,0.9}{
            \makecell[{{p{\linewidth}}}]{
                \texttt{\tiny{[GM$|$GM]}}
                \texttt{MUOVERE: W, (3, 2)} \\
            }
        }
    }
    & & \\ \\

    \theutterance \stepcounter{utterance}  
    & & & \multicolumn{2}{p{0.3\linewidth}}{
        \cellcolor[rgb]{0.9,0.9,0.9}{
            \makecell[{{p{\linewidth}}}]{
                \texttt{\tiny{[GM$|$GM]}}
                \texttt{W moved to (3, 2)} \\
            }
        }
    }
    & & \\ \\

    \theutterance \stepcounter{utterance}  
    & & \multicolumn{4}{p{0.6\linewidth}}{
        \cellcolor[rgb]{0.9,0.9,0.9}{
            \makecell[{{p{\linewidth}}}]{
                \texttt{\tiny{[P1$\langle$GM]}}
                \texttt{Spostato 'W' in (3,2) con successo. La tua griglia aggiornata appare così:} \\
\texttt{```} \\
\texttt{ 1234567} \\
\texttt{╔═══╤═╤═╗} \\
\texttt{║◌◌I│◌│◌║ 1} \\
\texttt{║◌◌W└─┤◌║ 2} \\
\texttt{║◌◌◌TC├─╢ 3} \\
\texttt{╟─┐◌┌─┤◌║ 4} \\
\texttt{║◌│◌│◌│◌║ 5} \\
\texttt{╟─┼─┼─┼─╢ 6} \\
\texttt{║◌│◌│H│◌║ 7} \\
\texttt{╚═╧═╧═╧═╝} \\
\texttt{```} \\
\texttt{Stai giocando il round 3 di un massimo di 20.} \\
\texttt{Avete accumulato collettivamente 2 di 12 penalità.} \\
\texttt{L'altro giocatore ha spostato un oggetto sulla sua griglia.} \\
\\ 
\\ 
\texttt{Qual è il tuo prossimo comando?} \\
            }
        }
    }
    & \\ \\

    \theutterance \stepcounter{utterance}  
    & \multicolumn{4}{p{0.6\linewidth}}{
        \cellcolor[rgb]{0.9,0.9,0.9}{
            \makecell[{{p{\linewidth}}}]{
                \texttt{\tiny{[P1$\rangle$GM]}}
                \texttt{MUOVERE: I, (3, 3)} \\
            }
        }
    }
    & & \\ \\

    \theutterance \stepcounter{utterance}  
    & & & \multicolumn{2}{p{0.3\linewidth}}{
        \cellcolor[rgb]{0.9,0.9,0.9}{
            \makecell[{{p{\linewidth}}}]{
                \texttt{\tiny{[GM$|$GM]}}
                \texttt{MUOVERE: I, (3, 3)} \\
            }
        }
    }
    & & \\ \\

    \theutterance \stepcounter{utterance}  
    & & & \multicolumn{2}{p{0.3\linewidth}}{
        \cellcolor[rgb]{0.9,0.9,0.9}{
            \makecell[{{p{\linewidth}}}]{
                \texttt{\tiny{[GM$|$GM]}}
                \texttt{I moved to (3, 3)} \\
            }
        }
    }
    & & \\ \\

    \theutterance \stepcounter{utterance}  
    & & \multicolumn{4}{p{0.6\linewidth}}{
        \cellcolor[rgb]{0.9,0.9,0.9}{
            \makecell[{{p{\linewidth}}}]{
                \texttt{\tiny{[GM$\rangle$P2]}}
                \texttt{Spostato 'W' in (3,2) con successo. La tua griglia aggiornata appare così:} \\
\texttt{```} \\
\texttt{ 1234567} \\
\texttt{╔═══╤═╤═╗} \\
\texttt{║◌◌◌│◌│◌║ 1} \\
\texttt{║◌◌W└─┤C║ 2} \\
\texttt{║T◌◌◌◌├─╢ 3} \\
\texttt{╟─┐◌┌─┤◌║ 4} \\
\texttt{║◌│◌│H│I║ 5} \\
\texttt{╟─┼─┼─┼─╢ 6} \\
\texttt{║◌│◌│◌│◌║ 7} \\
\texttt{╚═╧═╧═╧═╝} \\
\texttt{```} \\
\texttt{Stai giocando il round 4 di un massimo di 20.} \\
\texttt{Avete accumulato collettivamente 2 di 12 penalità.} \\
\texttt{L'altro giocatore ha spostato un oggetto sulla sua griglia.} \\
\\ 
\\ 
\texttt{Qual è il tuo prossimo comando?} \\
            }
        }
    }
    & \\ \\

    \theutterance \stepcounter{utterance}  
    & & & \multicolumn{4}{p{0.6\linewidth}}{
        \cellcolor[rgb]{0.9,0.9,0.9}{
            \makecell[{{p{\linewidth}}}]{
                \texttt{\tiny{[GM$\langle$P2]}}
                \texttt{MUOVERE: I, (3, 3)} \\
            }
        }
    }
     \\ \\

    \theutterance \stepcounter{utterance}  
    & & & \multicolumn{2}{p{0.3\linewidth}}{
        \cellcolor[rgb]{0.9,0.9,0.9}{
            \makecell[{{p{\linewidth}}}]{
                \texttt{\tiny{[GM$|$GM]}}
                \texttt{MUOVERE: I, (3, 3)} \\
            }
        }
    }
    & & \\ \\

    \theutterance \stepcounter{utterance}  
    & & & \multicolumn{2}{p{0.3\linewidth}}{
        \cellcolor[rgb]{0.9,0.9,0.9}{
            \makecell[{{p{\linewidth}}}]{
                \texttt{\tiny{[GM$|$GM]}}
                \texttt{I moved to (3, 3)} \\
            }
        }
    }
    & & \\ \\

    \theutterance \stepcounter{utterance}  
    & & \multicolumn{4}{p{0.6\linewidth}}{
        \cellcolor[rgb]{0.9,0.9,0.9}{
            \makecell[{{p{\linewidth}}}]{
                \texttt{\tiny{[P1$\langle$GM]}}
                \texttt{Spostato 'I' in (3,3) con successo. La tua griglia aggiornata appare così:} \\
\texttt{```} \\
\texttt{ 1234567} \\
\texttt{╔═══╤═╤═╗} \\
\texttt{║◌◌◌│◌│◌║ 1} \\
\texttt{║◌◌W└─┤◌║ 2} \\
\texttt{║◌◌ITC├─╢ 3} \\
\texttt{╟─┐◌┌─┤◌║ 4} \\
\texttt{║◌│◌│◌│◌║ 5} \\
\texttt{╟─┼─┼─┼─╢ 6} \\
\texttt{║◌│◌│H│◌║ 7} \\
\texttt{╚═╧═╧═╧═╝} \\
\texttt{```} \\
\texttt{Stai giocando il round 4 di un massimo di 20.} \\
\texttt{Avete accumulato collettivamente 2 di 12 penalità.} \\
\texttt{L'altro giocatore ha spostato un oggetto sulla sua griglia.} \\
\\ 
\\ 
\texttt{Qual è il tuo prossimo comando?} \\
            }
        }
    }
    & \\ \\

    \theutterance \stepcounter{utterance}  
    & \multicolumn{4}{p{0.6\linewidth}}{
        \cellcolor[rgb]{0.9,0.9,0.9}{
            \makecell[{{p{\linewidth}}}]{
                \texttt{\tiny{[P1$\rangle$GM]}}
                \texttt{MUOVERE: T, (3, 4)} \\
            }
        }
    }
    & & \\ \\

    \theutterance \stepcounter{utterance}  
    & & & \multicolumn{2}{p{0.3\linewidth}}{
        \cellcolor[rgb]{0.9,0.9,0.9}{
            \makecell[{{p{\linewidth}}}]{
                \texttt{\tiny{[GM$|$GM]}}
                \texttt{MUOVERE: T, (3, 4)} \\
            }
        }
    }
    & & \\ \\

    \theutterance \stepcounter{utterance}  
    & & & \multicolumn{2}{p{0.3\linewidth}}{
        \cellcolor[rgb]{0.9,0.9,0.9}{
            \makecell[{{p{\linewidth}}}]{
                \texttt{\tiny{[GM$|$GM]}}
                \texttt{T moved to (3, 4)} \\
            }
        }
    }
    & & \\ \\

    \theutterance \stepcounter{utterance}  
    & & \multicolumn{4}{p{0.6\linewidth}}{
        \cellcolor[rgb]{0.9,0.9,0.9}{
            \makecell[{{p{\linewidth}}}]{
                \texttt{\tiny{[GM$\rangle$P2]}}
                \texttt{Spostato 'I' in (3,3) con successo. La tua griglia aggiornata appare così:} \\
\texttt{```} \\
\texttt{ 1234567} \\
\texttt{╔═══╤═╤═╗} \\
\texttt{║◌◌◌│◌│◌║ 1} \\
\texttt{║◌◌W└─┤C║ 2} \\
\texttt{║T◌I◌◌├─╢ 3} \\
\texttt{╟─┐◌┌─┤◌║ 4} \\
\texttt{║◌│◌│H│◌║ 5} \\
\texttt{╟─┼─┼─┼─╢ 6} \\
\texttt{║◌│◌│◌│◌║ 7} \\
\texttt{╚═╧═╧═╧═╝} \\
\texttt{```} \\
\texttt{Stai giocando il round 5 di un massimo di 20.} \\
\texttt{Avete accumulato collettivamente 2 di 12 penalità.} \\
\texttt{L'altro giocatore ha spostato un oggetto sulla sua griglia.} \\
\\ 
\\ 
\texttt{Qual è il tuo prossimo comando?} \\
            }
        }
    }
    & \\ \\

    \theutterance \stepcounter{utterance}  
    & & & \multicolumn{4}{p{0.6\linewidth}}{
        \cellcolor[rgb]{0.9,0.9,0.9}{
            \makecell[{{p{\linewidth}}}]{
                \texttt{\tiny{[GM$\langle$P2]}}
                \texttt{MUOVERE: T, (3, 4)} \\
            }
        }
    }
     \\ \\

    \theutterance \stepcounter{utterance}  
    & & & \multicolumn{2}{p{0.3\linewidth}}{
        \cellcolor[rgb]{0.9,0.9,0.9}{
            \makecell[{{p{\linewidth}}}]{
                \texttt{\tiny{[GM$|$GM]}}
                \texttt{MUOVERE: T, (3, 4)} \\
            }
        }
    }
    & & \\ \\

    \theutterance \stepcounter{utterance}  
    & & & \multicolumn{2}{p{0.3\linewidth}}{
        \cellcolor[rgb]{0.9,0.9,0.9}{
            \makecell[{{p{\linewidth}}}]{
                \texttt{\tiny{[GM$|$GM]}}
                \texttt{T moved to (3, 4)} \\
            }
        }
    }
    & & \\ \\

    \theutterance \stepcounter{utterance}  
    & & \multicolumn{4}{p{0.6\linewidth}}{
        \cellcolor[rgb]{0.9,0.9,0.9}{
            \makecell[{{p{\linewidth}}}]{
                \texttt{\tiny{[P1$\langle$GM]}}
                \texttt{Spostato 'T' in (3,4) con successo. La tua griglia aggiornata appare così:} \\
\texttt{```} \\
\texttt{ 1234567} \\
\texttt{╔═══╤═╤═╗} \\
\texttt{║◌◌◌│◌│◌║ 1} \\
\texttt{║◌◌W└─┤◌║ 2} \\
\texttt{║◌◌I◌C├─╢ 3} \\
\texttt{╟─┐T┌─┤◌║ 4} \\
\texttt{║◌│◌│◌│◌║ 5} \\
\texttt{╟─┼─┼─┼─╢ 6} \\
\texttt{║◌│◌│H│◌║ 7} \\
\texttt{╚═╧═╧═╧═╝} \\
\texttt{```} \\
\texttt{Stai giocando il round 5 di un massimo di 20.} \\
\texttt{Avete accumulato collettivamente 2 di 12 penalità.} \\
\texttt{L'altro giocatore ha spostato un oggetto sulla sua griglia.} \\
\\ 
\\ 
\texttt{Qual è il tuo prossimo comando?} \\
            }
        }
    }
    & \\ \\

    \theutterance \stepcounter{utterance}  
    & \multicolumn{4}{p{0.6\linewidth}}{
        \cellcolor[rgb]{0.9,0.9,0.9}{
            \makecell[{{p{\linewidth}}}]{
                \texttt{\tiny{[P1$\rangle$GM]}}
                \texttt{MUOVERE: C, (3, 5)} \\
            }
        }
    }
    & & \\ \\

    \theutterance \stepcounter{utterance}  
    & & & \multicolumn{2}{p{0.3\linewidth}}{
        \cellcolor[rgb]{0.9,0.9,0.9}{
            \makecell[{{p{\linewidth}}}]{
                \texttt{\tiny{[GM$|$GM]}}
                \texttt{MUOVERE: C, (3, 5)} \\
            }
        }
    }
    & & \\ \\

    \theutterance \stepcounter{utterance}  
    & & & \multicolumn{2}{p{0.3\linewidth}}{
        \cellcolor[rgb]{0.9,0.9,0.9}{
            \makecell[{{p{\linewidth}}}]{
                \texttt{\tiny{[GM$|$GM]}}
                \texttt{C moved to (3, 5)} \\
            }
        }
    }
    & & \\ \\

    \theutterance \stepcounter{utterance}  
    & & \multicolumn{4}{p{0.6\linewidth}}{
        \cellcolor[rgb]{0.9,0.9,0.9}{
            \makecell[{{p{\linewidth}}}]{
                \texttt{\tiny{[GM$\rangle$P2]}}
                \texttt{Spostato 'T' in (3,4) con successo. La tua griglia aggiornata appare così:} \\
\texttt{```} \\
\texttt{ 1234567} \\
\texttt{╔═══╤═╤═╗} \\
\texttt{║◌◌◌│◌│◌║ 1} \\
\texttt{║◌◌W└─┤C║ 2} \\
\texttt{║◌◌I◌◌├─╢ 3} \\
\texttt{╟─┐T┌─┤◌║ 4} \\
\texttt{║◌│◌│H│◌║ 5} \\
\texttt{╟─┼─┼─┼─╢ 6} \\
\texttt{║◌│◌│◌│◌║ 7} \\
\texttt{╚═╧═╧═╧═╝} \\
\texttt{```} \\
\texttt{Stai giocando il round 6 di un massimo di 20.} \\
\texttt{Avete accumulato collettivamente 2 di 12 penalità.} \\
\texttt{L'altro giocatore ha spostato un oggetto sulla sua griglia.} \\
\\ 
\\ 
\texttt{Qual è il tuo prossimo comando?} \\
            }
        }
    }
    & \\ \\

    \theutterance \stepcounter{utterance}  
    & & & \multicolumn{4}{p{0.6\linewidth}}{
        \cellcolor[rgb]{0.9,0.9,0.9}{
            \makecell[{{p{\linewidth}}}]{
                \texttt{\tiny{[GM$\langle$P2]}}
                \texttt{MUOVERE: C, (3, 5)} \\
            }
        }
    }
     \\ \\

    \theutterance \stepcounter{utterance}  
    & & & \multicolumn{2}{p{0.3\linewidth}}{
        \cellcolor[rgb]{0.9,0.9,0.9}{
            \makecell[{{p{\linewidth}}}]{
                \texttt{\tiny{[GM$|$GM]}}
                \texttt{MUOVERE: C, (3, 5)} \\
            }
        }
    }
    & & \\ \\

    \theutterance \stepcounter{utterance}  
    & & & \multicolumn{2}{p{0.3\linewidth}}{
        \cellcolor[rgb]{0.9,0.9,0.9}{
            \makecell[{{p{\linewidth}}}]{
                \texttt{\tiny{[GM$|$GM]}}
                \texttt{C moved to (3, 5)} \\
            }
        }
    }
    & & \\ \\

    \theutterance \stepcounter{utterance}  
    & & \multicolumn{4}{p{0.6\linewidth}}{
        \cellcolor[rgb]{0.9,0.9,0.9}{
            \makecell[{{p{\linewidth}}}]{
                \texttt{\tiny{[P1$\langle$GM]}}
                \texttt{Spostato 'C' in (3,5) con successo. La tua griglia aggiornata appare così:} \\
\texttt{```} \\
\texttt{ 1234567} \\
\texttt{╔═══╤═╤═╗} \\
\texttt{║◌◌◌│◌│◌║ 1} \\
\texttt{║◌◌W└─┤◌║ 2} \\
\texttt{║◌◌I◌◌├─╢ 3} \\
\texttt{╟─┐T┌─┤◌║ 4} \\
\texttt{║◌│C│◌│◌║ 5} \\
\texttt{╟─┼─┼─┼─╢ 6} \\
\texttt{║◌│◌│H│◌║ 7} \\
\texttt{╚═╧═╧═╧═╝} \\
\texttt{```} \\
\texttt{Stai giocando il round 6 di un massimo di 20.} \\
\texttt{Avete accumulato collettivamente 2 di 12 penalità.} \\
\texttt{L'altro giocatore ha spostato un oggetto sulla sua griglia.} \\
\\ 
\\ 
\texttt{Qual è il tuo prossimo comando?} \\
            }
        }
    }
    & \\ \\

    \theutterance \stepcounter{utterance}  
    & \multicolumn{4}{p{0.6\linewidth}}{
        \cellcolor[rgb]{0.9,0.9,0.9}{
            \makecell[{{p{\linewidth}}}]{
                \texttt{\tiny{[P1$\rangle$GM]}}
                \texttt{MUOVERE: H, (3, 7)} \\
            }
        }
    }
    & & \\ \\

    \theutterance \stepcounter{utterance}  
    & & & \multicolumn{2}{p{0.3\linewidth}}{
        \cellcolor[rgb]{0.9,0.9,0.9}{
            \makecell[{{p{\linewidth}}}]{
                \texttt{\tiny{[GM$|$GM]}}
                \texttt{MUOVERE: H, (3, 7)} \\
            }
        }
    }
    & & \\ \\

    \theutterance \stepcounter{utterance}  
    & & & \multicolumn{2}{p{0.3\linewidth}}{
        \cellcolor[rgb]{0.9,0.9,0.9}{
            \makecell[{{p{\linewidth}}}]{
                \texttt{\tiny{[GM$|$GM]}}
                \texttt{H moved to (3, 7)} \\
            }
        }
    }
    & & \\ \\

    \theutterance \stepcounter{utterance}  
    & & \multicolumn{4}{p{0.6\linewidth}}{
        \cellcolor[rgb]{0.9,0.9,0.9}{
            \makecell[{{p{\linewidth}}}]{
                \texttt{\tiny{[GM$\rangle$P2]}}
                \texttt{Spostato 'C' in (3,5) con successo. La tua griglia aggiornata appare così:} \\
\texttt{```} \\
\texttt{ 1234567} \\
\texttt{╔═══╤═╤═╗} \\
\texttt{║◌◌◌│◌│◌║ 1} \\
\texttt{║◌◌W└─┤◌║ 2} \\
\texttt{║◌◌I◌◌├─╢ 3} \\
\texttt{╟─┐T┌─┤◌║ 4} \\
\texttt{║◌│C│H│◌║ 5} \\
\texttt{╟─┼─┼─┼─╢ 6} \\
\texttt{║◌│◌│◌│◌║ 7} \\
\texttt{╚═╧═╧═╧═╝} \\
\texttt{```} \\
\texttt{Stai giocando il round 7 di un massimo di 20.} \\
\texttt{Avete accumulato collettivamente 2 di 12 penalità.} \\
\texttt{L'altro giocatore ha spostato un oggetto sulla sua griglia.} \\
\\ 
\\ 
\texttt{Qual è il tuo prossimo comando?} \\
            }
        }
    }
    & \\ \\

    \theutterance \stepcounter{utterance}  
    & & & \multicolumn{4}{p{0.6\linewidth}}{
        \cellcolor[rgb]{0.9,0.9,0.9}{
            \makecell[{{p{\linewidth}}}]{
                \texttt{\tiny{[GM$\langle$P2]}}
                \texttt{MUOVERE: H, (3, 7)} \\
            }
        }
    }
     \\ \\

    \theutterance \stepcounter{utterance}  
    & & & \multicolumn{2}{p{0.3\linewidth}}{
        \cellcolor[rgb]{0.9,0.9,0.9}{
            \makecell[{{p{\linewidth}}}]{
                \texttt{\tiny{[GM$|$GM]}}
                \texttt{MUOVERE: H, (3, 7)} \\
            }
        }
    }
    & & \\ \\

    \theutterance \stepcounter{utterance}  
    & & & \multicolumn{2}{p{0.3\linewidth}}{
        \cellcolor[rgb]{0.9,0.9,0.9}{
            \makecell[{{p{\linewidth}}}]{
                \texttt{\tiny{[GM$|$GM]}}
                \texttt{H moved to (3, 7)} \\
            }
        }
    }
    & & \\ \\

    \theutterance \stepcounter{utterance}  
    & & \multicolumn{4}{p{0.6\linewidth}}{
        \cellcolor[rgb]{0.9,0.9,0.9}{
            \makecell[{{p{\linewidth}}}]{
                \texttt{\tiny{[P1$\langle$GM]}}
                \texttt{Spostato 'H' in (3,7) con successo. La tua griglia aggiornata appare così:} \\
\texttt{```} \\
\texttt{ 1234567} \\
\texttt{╔═══╤═╤═╗} \\
\texttt{║◌◌◌│◌│◌║ 1} \\
\texttt{║◌◌W└─┤◌║ 2} \\
\texttt{║◌◌I◌◌├─╢ 3} \\
\texttt{╟─┐T┌─┤◌║ 4} \\
\texttt{║◌│C│◌│◌║ 5} \\
\texttt{╟─┼─┼─┼─╢ 6} \\
\texttt{║◌│H│◌│◌║ 7} \\
\texttt{╚═╧═╧═╧═╝} \\
\texttt{```} \\
\texttt{Stai giocando il round 7 di un massimo di 20.} \\
\texttt{Avete accumulato collettivamente 2 di 12 penalità.} \\
\texttt{L'altro giocatore ha spostato un oggetto sulla sua griglia.} \\
\\ 
\\ 
\texttt{Qual è il tuo prossimo comando?} \\
            }
        }
    }
    & \\ \\

    \theutterance \stepcounter{utterance}  
    & \multicolumn{4}{p{0.6\linewidth}}{
        \cellcolor[rgb]{0.9,0.9,0.9}{
            \makecell[{{p{\linewidth}}}]{
                \texttt{\tiny{[P1$\rangle$GM]}}
                \texttt{DIRE: finito?} \\
            }
        }
    }
    & & \\ \\

    \theutterance \stepcounter{utterance}  
    & & & \multicolumn{2}{p{0.3\linewidth}}{
        \cellcolor[rgb]{0.9,0.9,0.9}{
            \makecell[{{p{\linewidth}}}]{
                \texttt{\tiny{[GM$|$GM]}}
                \texttt{DIRE: finito?} \\
            }
        }
    }
    & & \\ \\

    \theutterance \stepcounter{utterance}  
    & & \multicolumn{4}{p{0.6\linewidth}}{
        \cellcolor[rgb]{0.9,0.9,0.9}{
            \makecell[{{p{\linewidth}}}]{
                \texttt{\tiny{[GM$\rangle$P2]}}
                \texttt{Spostato 'H' in (3,7) con successo. La tua griglia aggiornata appare così:} \\
\texttt{```} \\
\texttt{ 1234567} \\
\texttt{╔═══╤═╤═╗} \\
\texttt{║◌◌◌│◌│◌║ 1} \\
\texttt{║◌◌W└─┤◌║ 2} \\
\texttt{║◌◌I◌◌├─╢ 3} \\
\texttt{╟─┐T┌─┤◌║ 4} \\
\texttt{║◌│C│◌│◌║ 5} \\
\texttt{╟─┼─┼─┼─╢ 6} \\
\texttt{║◌│H│◌│◌║ 7} \\
\texttt{╚═╧═╧═╧═╝} \\
\texttt{```} \\
\texttt{Stai giocando il round 8 di un massimo di 20.} \\
\texttt{Avete accumulato collettivamente 2 di 12 penalità.} \\
\texttt{L'altro giocatore ha inviato questo messaggio:} \\
\\ 
\\ 
\texttt{"finito?"} \\
\\ 
\\ 
\texttt{Qual è il tuo prossimo comando?} \\
            }
        }
    }
    & \\ \\

    \theutterance \stepcounter{utterance}  
    & & & \multicolumn{4}{p{0.6\linewidth}}{
        \cellcolor[rgb]{0.9,0.9,0.9}{
            \makecell[{{p{\linewidth}}}]{
                \texttt{\tiny{[GM$\langle$P2]}}
                \texttt{DIRE: finito!} \\
            }
        }
    }
     \\ \\

    \theutterance \stepcounter{utterance}  
    & & & \multicolumn{2}{p{0.3\linewidth}}{
        \cellcolor[rgb]{0.9,0.9,0.9}{
            \makecell[{{p{\linewidth}}}]{
                \texttt{\tiny{[GM$|$GM]}}
                \texttt{DIRE: finito!} \\
            }
        }
    }
    & & \\ \\

    \theutterance \stepcounter{utterance}  
    & & & \multicolumn{2}{p{0.3\linewidth}}{
        \cellcolor[rgb]{0.9,0.9,0.9}{
            \makecell[{{p{\linewidth}}}]{
                \texttt{\tiny{[GM$|$GM]}}
                \texttt{true} \\
            }
        }
    }
    & & \\ \\

    \theutterance \stepcounter{utterance}  
    & & & \multicolumn{2}{p{0.3\linewidth}}{
        \cellcolor[rgb]{0.9,0.9,0.9}{
            \makecell[{{p{\linewidth}}}]{
                \texttt{\tiny{[GM$|$GM]}}
                \texttt{Initial states:} \\
\texttt{```} \\
\texttt{Player 1:} \\
\texttt{ 1234567} \\
\texttt{╔═══╤═╤═╗} \\
\texttt{║◌◌I│W│◌║ 1} \\
\texttt{║◌◌◌└─┤◌║ 2} \\
\texttt{║◌◌◌TC├─╢ 3} \\
\texttt{╟─┐◌┌─┤◌║ 4} \\
\texttt{║◌│◌│◌│◌║ 5} \\
\texttt{╟─┼─┼─┼─╢ 6} \\
\texttt{║◌│◌│H│◌║ 7} \\
\texttt{╚═╧═╧═╧═╝} \\
\\ 
\\ 
\texttt{Player 2:} \\
\texttt{ 1234567} \\
\texttt{╔═══╤═╤═╗} \\
\texttt{║◌◌◌│◌│◌║ 1} \\
\texttt{║◌W◌└─┤C║ 2} \\
\texttt{║T◌◌◌◌├─╢ 3} \\
\texttt{╟─┐◌┌─┤◌║ 4} \\
\texttt{║◌│◌│H│I║ 5} \\
\texttt{╟─┼─┼─┼─╢ 6} \\
\texttt{║◌│◌│◌│◌║ 7} \\
\texttt{╚═╧═╧═╧═╝} \\
\\ 
\texttt{```} \\
            }
        }
    }
    & & \\ \\

    \theutterance \stepcounter{utterance}  
    & & & \multicolumn{2}{p{0.3\linewidth}}{
        \cellcolor[rgb]{0.9,0.9,0.9}{
            \makecell[{{p{\linewidth}}}]{
                \texttt{\tiny{[GM$|$GM]}}
                \texttt{End states:} \\
\texttt{```} \\
\texttt{Player 1:} \\
\texttt{ 1234567} \\
\texttt{╔═══╤═╤═╗} \\
\texttt{║◌◌◌│◌│◌║ 1} \\
\texttt{║◌◌W└─┤◌║ 2} \\
\texttt{║◌◌I◌◌├─╢ 3} \\
\texttt{╟─┐T┌─┤◌║ 4} \\
\texttt{║◌│C│◌│◌║ 5} \\
\texttt{╟─┼─┼─┼─╢ 6} \\
\texttt{║◌│H│◌│◌║ 7} \\
\texttt{╚═╧═╧═╧═╝} \\
\\ 
\\ 
\texttt{Player 2:} \\
\texttt{ 1234567} \\
\texttt{╔═══╤═╤═╗} \\
\texttt{║◌◌◌│◌│◌║ 1} \\
\texttt{║◌◌W└─┤◌║ 2} \\
\texttt{║◌◌I◌◌├─╢ 3} \\
\texttt{╟─┐T┌─┤◌║ 4} \\
\texttt{║◌│C│◌│◌║ 5} \\
\texttt{╟─┼─┼─┼─╢ 6} \\
\texttt{║◌│H│◌│◌║ 7} \\
\texttt{╚═╧═╧═╧═╝} \\
\\ 
\texttt{```} \\
            }
        }
    }
    & & \\ \\

    \theutterance \stepcounter{utterance}  
    & & & \multicolumn{2}{p{0.3\linewidth}}{
        \cellcolor[rgb]{0.9,0.9,0.9}{
            \makecell[{{p{\linewidth}}}]{
                \texttt{\tiny{[GM$|$GM]}}
                \texttt{* success: True} \\
\texttt{* lose: False} \\
\texttt{* aborted: False} \\
\texttt{{-}{-}{-}{-}{-}{-}{-}} \\
\texttt{* Shifts: 4.00} \\
\texttt{* Max Shifts: 8.00} \\
\texttt{* Min Shifts: 4.00} \\
\texttt{* End Distance Sum: 0.00} \\
\texttt{* Init Distance Sum: 16.06} \\
\texttt{* Expected Distance Sum: 20.95} \\
\texttt{* Penalties: 2.00} \\
\texttt{* Max Penalties: 12.00} \\
\texttt{* Rounds: 8.00} \\
\texttt{* Max Rounds: 20.00} \\
\texttt{* Object Count: 5.00} \\
            }
        }
    }
    & & \\ \\

\end{supertabular}
}

\end{document}
